% Cover Letter — IEEE Transactions on Medical Imaging submission
% Compile: pdflatex mic-cover-letter-ieee-tmi.tex
\documentclass[12pt]{article}

\usepackage[margin=1in]{geometry}
\usepackage{xcolor}
\usepackage{hyperref}
\usepackage{parskip}
\usepackage{microtype}

\hypersetup{
  colorlinks=true,
  urlcolor=blue!70!black,
}

\pagestyle{empty}

\begin{document}

\begin{flushright}
Kuldeep Singh\\
Innovaccer Inc.\\
\href{mailto:kuldeep.singh@innovaccer.com}{kuldeep.singh@innovaccer.com}\\
ORCID: 0009-0004-8476-0118\\[6pt]
May 2, 2026
\end{flushright}

\bigskip

Editor-in-Chief\\
\textit{IEEE Transactions on Medical Imaging}

\bigskip

\noindent Dear Editor-in-Chief,

\medskip

I am submitting the manuscript \textbf{``A 16-Bit-Native Entropy Coding
Pipeline for Lossless Medical Image Compression''} for consideration as a Full
Paper in \textit{IEEE Transactions on Medical Imaging}.

\medskip
\noindent\textbf{Summary.}
This paper presents MIC, a lossless compression codec for high-bit-depth
Digital Imaging and Communications in Medicine (DICOM) images. The central
contribution is a pipeline that performs spatial prediction, 16-bit-native
run-length encoding, and large-alphabet Finite State Entropy (FSE) coding
directly on 16-bit residuals, rather than splitting pixels into byte pairs
before entropy coding. I show theoretically and empirically that byte-splitting
incurs a structural information loss of 0.3--1.1 bits per symbol on
delta-coded medical residuals (8--22\% of total entropy), and that a
16-bit-native coder recovers this cost. The paper additionally contributes a
four-state interleaved FSE decoder that improves decompression throughput by a
geometric mean of 68\% on ARM64 and 43\% on AMD64 with no change to
compressed size, a parallel strip (PICS) mode that exceeds High-Throughput
JPEG~2000 (HTJ2K) decompression speed on all 21 tested images on both
platforms, and a approximately 20-kilobyte pure JavaScript decoder enabling
client-side in-browser DICOM viewing without server-side transcoding.

\medskip
\noindent\textbf{Relevance to IEEE TMI.}
The manuscript addresses lossless compression for clinical DICOM data---a
core challenge in medical image storage, PACS transmission, and diagnostic
viewing of 10--16 bit pixel data. Results are reported on 21 de-identified
clinical DICOM images spanning nine modalities (CT, MRI, mammography,
computed radiography, radiography, fluoroscopy, nuclear medicine, secondary
capture, X-ray), with fair in-process comparison against HTJ2K (OpenJPH)
and JPEG-LS (CharLS). The combination of compression-ratio analysis,
throughput benchmarking, and browser deployment feasibility is directly
relevant to the TMI readership working on medical image informatics systems.

\medskip
\noindent\textbf{Originality.}
This work has not been submitted to, or published in, any other journal.
The manuscript represents a complete original work; no prior conference or
workshop version has been published.

\medskip
\noindent\textbf{Open-source reproducibility.}
The complete implementation, benchmark harness, and test data references are
available at \url{https://github.com/pappuks/medical-image-codec}. All results
are reproducible using the commands documented in the repository.

\medskip
\noindent\textbf{Ethics.}
The evaluation dataset consists entirely of publicly available, de-identified
DICOM images from the NEMA WG-04 compsamples archive, the NEMA 1997~CD, and
the Clunie Breast Tomosynthesis public case. No human subjects research was
conducted; no institutional ethics approval is required per DICOM PS~3.15
Appendix~E.

\medskip
\noindent\textbf{Conflict of interest.}
The author declares no competing financial interests or personal relationships
that could have influenced this work.

\medskip
\noindent\textbf{Suggested Associate Editor area.}
Medical Image Compression; Medical Image Storage and Transmission; Medical
Imaging Informatics Systems.

\medskip
\noindent\textbf{Note on suggested reviewers.}
I note that the paper compares against OpenJPH and CharLS; authors of those
libraries may have a conflict of interest and could reasonably be excluded
at the Editor's discretion.

\bigskip

I confirm this is a single-author submission and that the manuscript has been
prepared in accordance with IEEE Transactions on Medical Imaging author
guidelines, including the 10-page limit, the 250-word abstract limit, and the
use of de-identified public datasets.

\bigskip

\noindent Sincerely,

\bigskip\bigskip

\noindent Kuldeep Singh\\
Innovaccer Inc.\\
\href{mailto:kuldeep.singh@innovaccer.com}{kuldeep.singh@innovaccer.com}\\
ORCID: 0009-0004-8476-0118\\
\url{https://github.com/pappuks/medical-image-codec}

\end{document}
